\section{Cronograma y Presupuesto}\label{sec:cronograma}

\subsection{Cronograma}

El anteproyecto estimó un período de ejecución de aproximadamente diez semanas a partir de la aprobación, durante las cuales se desarrollarían la construcción y validación de la taxonomía, la selección de repositorios, la implementación del analizador estático, la extracción de métricas y el análisis estadístico. Esa estimación describe la intención del diseño, pero no permite reconstruir un cronograma detallado actividad por actividad: los registros del proyecto no conservan fechas de inicio y fin para cada una de esas actividades, ni la fecha formal de aprobación que serviría de ancla.

Lo que sí puede documentarse con evidencia son las marcas temporales de la ejecución técnica. La extracción de metadatos de la muestra se registró el 23 de noviembre de 2026 a las 20:13:30 UTC. La verificación autenticada de los repositorios preseleccionados se registró el mismo día a las 21:10:48 UTC. El escaneo de los 23 repositorios se ejecutó entre las 21:56:40 y las 21:59:41 UTC, y las etapas posteriores de métricas, análisis estadístico y generación de tablas y figuras se ejecutaron a continuación. Estos registros permiten afirmar que la fase de recolección y análisis automatizado se completó en una ventana de aproximadamente dos horas, pero no permiten fechar el diseño, la revisión de literatura, la escritura ni la corrección del informe.

Por esa razón, este documento no incluye un diagrama de Gantt. Un Gantt exige fechas de inicio y fin por actividad, hitos y dependencias; construirlo con valores estimados o retrocedidos introduciría información que no está respaldada por los registros y contradiría el principio de integridad científica que orienta el informe. La omisión es deliberada y se declara como una limitación administrativa, no como un olvido. Si el programa académico exige el diagrama, este deberá elaborarse a partir de la bitácora real del investigador o de las actas de aprobación, que no forman parte de los artefactos disponibles.

\subsection{Presupuesto}

Tampoco existe un registro contable del proyecto. No hay facturas, cotizaciones, comprobantes de pago, contratos de personal ni asignaciones presupuestales que permitan reportar costos monetarios por rubro. Por tanto, no se presenta una tabla de presupuesto y no se atribuyen valores a horas de trabajo, equipos, licencias, servicios en la nube o desplazamientos. Hacerlo sería una invención.

Lo que puede describirse es el conjunto de recursos efectivamente utilizados, sin asignarles un valor monetario. El trabajo se apoyó en herramientas de código abierto y en infraestructura preexistente: Python y su ecosistema científico (pandas, SciPy, Matplotlib), Radon para métricas de código, la API pública de GitHub para la recolección y verificación, Git para el control de versiones y una distribución local de LaTeX (MiKTeX) para la compilación del documento. No se adquirieron licencias comerciales, no se contrató personal y no se utilizaron servicios de cómputo en la nube con costo registrado. El almacenamiento y el procesamiento se realizaron en el equipo del investigador.

En términos de esfuerzo, la ejecución técnica automatizada fue breve, pero estuvo precedida por el diseño de la taxonomía, la implementación y depuración de los detectores, la definición de las decisiones metodológicas y la revisión de la literatura, actividades cuya dedicación horaria no quedó registrada. El informe reconoce esa ausencia y no la sustituye por estimaciones. La consecuencia práctica es que el proyecto puede auditarse en cuanto a sus artefactos y resultados, pero no en cuanto a su costo o a su calendarización formal.
